\section{Anatomy of a starved visual channel: openVINS on the LEO rover}
\label{sec:ovstarvation}

Over a single instrumented session (2026-08-20) the visual--inertial
estimator was taken from ``it drifts'' to a quantified diagnosis. Three
defects were found, measured and corrected; seven candidate explanations for
the residual failure were eliminated by measurement rather than by argument.
This section records the mechanisms, because each is a trap that costs a night
to rediscover.

The single number that governs everything below is the count of visual
features the filter actually \emph{consumes} per update. openVINS is
configured with \texttt{max\_msckf\_in\_update: 40}. Measured on this robot,
over $1793$ updates while driving:
\begin{equation}
\bar{n}_{\mathrm{feat}} = 0.75,
\qquad
n_{\max} = 25,
\qquad
\frac{\bar{n}_{\mathrm{feat}}}{40} = 1.9\,\%.
\label{eq:featstarve}
\end{equation}
The estimator is therefore not a visual--inertial filter in practice but an
inertial integrator with occasional visual corrections. Every symptom follows
from that.

% ─────────────────────────────────────────────────────────────────────────
\subsection{The measurement funnel, and where it empties}
\label{sec:funnel}

The front end is \emph{not} idle. Its debug overlay draws roughly
$1.19\times10^{4}$ coloured pixels per frame, so points are detected and
tracked. The loss happens between detection and the state update.

\begin{figure}[htbp]
\centering
\begin{tikzpicture}[font=\small, >=Latex, scale=0.95]
  \tikzset{
    stage/.style={draw, rounded corners=1pt, minimum height=8mm,
                  align=center, font=\small},
    good/.style={stage, fill=accentblue!8, draw=accentblue!55},
    bad/.style={stage, fill=alertred!8, draw=alertred!60},
    lost/.style={font=\scriptsize\itshape, color=alertred!85, align=left}
  }
  \node[good, minimum width=62mm] (det)  at (0,0)      {FAST/KLT detection \& tracking
                                                        \;\textbf{$\sim$300 points}};
  \node[good, minimum width=50mm] (mat)  at (0,-1.5)   {stereo match, RANSAC epipolar
                                                        \;(\texttt{findFundamentalMat})};
  \node[bad,  minimum width=38mm] (tri)  at (0,-3.0)   {triangulation
                                                        \;(\texttt{FeatureInitializer})};
  \node[bad,  minimum width=26mm] (chi)  at (0,-4.5)   {$\chi^2$ gate};
  \node[bad,  minimum width=16mm] (upd)  at (0,-6.0)   {\textbf{0.75}};

  \draw[->, thick] (det) -- (mat);
  \draw[->, thick] (mat) -- (tri);
  \draw[->, thick] (tri) -- (chi);
  \draw[->, thick] (chi) -- (upd);

  \node[lost, anchor=west] at (3.6,-2.25) {epipolar constraint is estimated\\
                                           from the \emph{data}, not from\\
                                           the configured extrinsic};
  \node[lost, anchor=west] at (3.6,-3.75) {rejects if
                                           $\|\mathbf{p}_f\|/b_{\max} > 40$\\
                                           $\Rightarrow$ $d_{\max}=3.6$\,m at rest};
  \node[lost, anchor=west] at (3.6,-5.25) {loosening $\times 5$ changed\\
                                           $0.58 \to 0.75$: not the gate};
  \node[font=\scriptsize, anchor=east, color=alertred] at (-1.2,-6.0)
        {features per update};
\end{tikzpicture}
\caption[Feature funnel in openVINS]{Where the visual measurements are lost.
Detection is healthy; the collapse happens downstream. Each annotation is a
hypothesis that was tested and its measured outcome. The stereo matching stage
cannot be blamed on a mis-declared extrinsic, because \texttt{TrackKLT} derives
the epipolar constraint by RANSAC from the observed correspondences.}
\label{fig:featfunnel}
\end{figure}

% ─────────────────────────────────────────────────────────────────────────
\subsection{Defect 1: a circular ZUPT lock}
\label{sec:zuptlock}

Zero-velocity updates are gated by an estimated-velocity threshold
$v_{\max}$ (\texttt{zupt\_max\_velocity}). The update is applied only when
\begin{equation}
\|\hat{\mathbf{v}}\| < v_{\max}.
\label{eq:zuptgate}
\end{equation}
A residual specific force $\varepsilon$ integrates into the velocity estimate
between corrections,
\begin{equation}
\|\hat{\mathbf{v}}(t)\| \;\approx\; \varepsilon\,t ,
\label{eq:vgrowth}
\end{equation}
so once $\|\hat{\mathbf{v}}\|$ exceeds $v_{\max}$ the gate \eqref{eq:zuptgate}
can never re-open. The estimate is trapped: the quantity that decides whether
the correction runs is precisely the quantity the correction was meant to fix.

\begin{figure}[htbp]
\centering
\begin{tikzpicture}[font=\small, >=Latex, node distance=8mm and 20mm]
  \tikzset{blk/.style={draw, rounded corners=1pt, minimum height=8mm,
                       minimum width=30mm, align=center, font=\small}}
  \node[blk, fill=alertred!8, draw=alertred!60] (bias) {accel residual
                                                        $\varepsilon$};
  \node[blk, right=28mm of bias, fill=codebg]   (int)  {$\int\!\mathrm{d}t$};
  \node[blk, below=14mm of int, fill=codebg]    (vel)  {$\|\hat{\mathbf{v}}\|$};
  \node[blk, left=28mm of vel, fill=alertred!8, draw=alertred!60]
                                                (gate) {gate\\
                                                        $\|\hat{\mathbf{v}}\|<v_{\max}$?};
  \draw[->, thick] (bias) -- (int);
  \draw[->, thick] (int)  -- (vel);
  \draw[->, thick] (vel)  -- (gate);
  \draw[->, thick, alertred] (gate) -- node[left, font=\scriptsize, align=center]
        {\textsc{no}\\ no correction} (bias);
  \node[font=\scriptsize\itshape, align=center, color=alertred!85]
        at (0,-2.6) {the loop has no exit above $v_{\max}$};
\end{tikzpicture}
\caption[The ZUPT circular lock]{The lock. Above the threshold the negative
feedback path is severed, and the residual keeps driving the estimate.}
\label{fig:zuptlock}
\end{figure}

The prediction of \eqref{eq:vgrowth}--\eqref{eq:zuptgate} is that
$\|\hat{\mathbf{v}}\|$ settles slightly \emph{above} whatever
$v_{\max}$ is set to, for any value of $v_{\max}$. This was tested directly:

\begin{table}[htbp]
\centering
\begin{tabular}{@{}lrr@{}}
\toprule
$v_{\max}$ (m/s) & plateau of $\|\hat{\mathbf{v}}\|$ (m/s) & ratio \\
\midrule
$0.1$  & $1.295$ & --- \\
$1.5$  & $1.588$ & $1.06$ \\
\bottomrule
\end{tabular}
\caption[ZUPT plateau versus threshold]{The plateau tracks the threshold,
confirming \eqref{eq:zuptgate} as the binding constraint. At $v_{\max}=1.5$
the estimate converged monotonically to $1.588$\,m/s over $55$\,s and stayed
there, with the wheels reading exactly zero throughout.}
\label{tab:zuptplateau}
\end{table}

Raising the threshold only moves the plateau. The gate was therefore
neutralised ($v_{\max}=50$), leaving the two constraints that do not depend on
the filter's own possibly-wrong state: the image-disparity test
(\texttt{zupt\_max\_disparity}$\,=0.5$) and the $\chi^2$ test. Measured
outcome at rest: $\|\hat{\mathbf{v}}\|$ from $1.295$ to
$3\times10^{-3}$\,m/s, and cumulative path over $45$\,s from $95.4$\,m to
$0.29$\,m.

% ─────────────────────────────────────────────────────────────────────────
\subsection{Defect 2: a calibration constant that expires}
\label{sec:accelexpire}

The CORE2 accelerometer under-reads. The correction is a scalar applied to the
raw specific force,
\begin{equation}
\mathbf{a}_{\mathrm{corr}} = s\,\mathbf{a}_{\mathrm{raw}},
\qquad
s = \frac{g_{\mathrm{local}}}{\|\bar{\mathbf{a}}_{\mathrm{raw}}\|_{\text{at rest}}},
\label{eq:accelscaleov}
\end{equation}
which is legitimate because $\|\mathbf{a}\|$ at rest equals the local gravity
magnitude \emph{independently of orientation} --- a norm is rotation
invariant. That is the same argument that licenses periodic gyro-bias
re-estimation, and it is what makes \eqref{eq:accelscaleov} safe to re-run in
the field.

The constant was measured once, on 2026-07-28, and frozen. Three weeks later
the sensor had moved:

\begin{table}[htbp]
\centering
\begin{tabular}{@{}lrrr@{}}
\toprule
 & 2026-07-28 & 2026-08-20 & drift \\
\midrule
$\|\bar{\mathbf{a}}_{\mathrm{raw}}\|$ (m/s$^2$) & $8.7525$ & $8.6531$ & $-1.14\,\%$ \\
required $s$                                    & $1.1186$ & $1.1314$ & $+1.14\,\%$ \\
residual with frozen $s$ (m/s$^2$)              & $\approx 0$ & $-0.111$ & --- \\
\bottomrule
\end{tabular}
\caption[Accelerometer scale drift over three weeks]{A calibration constant
written into a launch file silently expires. Nothing fails; the estimator
simply diverges.}
\label{tab:accelexpire}
\end{table}

A constant specific-force error $\varepsilon$ integrates twice into position,
\begin{equation}
d(t) = \tfrac{1}{2}\,\varepsilon\,t^{2},
\label{eq:doubleintov}
\end{equation}
so the $-0.111$\,m/s$^2$ residual predicts
$d(600\,\mathrm{s}) = \tfrac{1}{2}(0.111)(600)^2 \approx 2.0\times10^{4}$\,m
--- twenty kilometres over a ten-minute run, which matches the observed
runaway in order of magnitude. Re-measuring gave a residual of
$+1.7\times10^{-3}$\,m/s$^2$, a reduction by a factor of $65$.

The lesson is structural rather than numerical: the fix is not the new value
but the mechanism that keeps it current. A periodic re-estimation
(\texttt{accel\_recal\_period}) now runs the same gated stationarity test used
for the gyro, so \eqref{eq:accelscaleov} is re-evaluated at confirmed stops.

% ─────────────────────────────────────────────────────────────────────────
\subsection{Defect 3: the triangulation baseline ratio}
\label{sec:baselineratio}

\texttt{ov\_core}'s feature initialiser rejects a triangulated point when
\begin{equation}
\frac{\|\mathbf{p}_f\|}{b_{\max}} > \beta,
\qquad \beta = \texttt{fi\_max\_baseline} = 40 \;\text{(default)},
\label{eq:baselinegate}
\end{equation}
where $b_{\max}$ is the largest baseline available among the sliding-window
clones. Rearranged, \eqref{eq:baselinegate} is a hard distance ceiling:
\begin{equation}
d_{\max} = \beta\, b_{\max}.
\label{eq:distceiling}
\end{equation}

\begin{figure}[htbp]
\centering
\begin{tikzpicture}[font=\small, >=Latex, scale=1.0]
  % deux centres optiques
  \fill (0,0) circle (1.6pt) node[below=2pt, font=\scriptsize] {$C_0$};
  \fill (1.6,0) circle (1.6pt) node[below=2pt, font=\scriptsize] {$C_1$};
  \draw[<->, accentblue, thick] (0,-0.45) -- node[below, font=\scriptsize,
        color=accentblue] {$b_{\max}$} (1.6,-0.45);
  % point lointain
  \fill[alertred] (5.4,2.5) circle (1.8pt) node[right=2pt, font=\scriptsize,
        color=alertred] {$\mathbf{p}_f$};
  \draw[black!60] (0,0) -- (5.4,2.5);
  \draw[black!60] (1.6,0) -- (5.4,2.5);
  % angle de parallaxe
  \draw[alertred, thick] (4.75,2.2) arc (200:222:1.4);
  \node[font=\scriptsize, color=alertred] at (4.35,2.62) {$\gamma$};
  % cone acceptable
  \draw[dashed, accentblue!70] (0.8,0) circle (2.6);
  \node[font=\scriptsize, color=accentblue, anchor=west] at (3.5,0.55)
        {$d_{\max}=\beta\,b_{\max}$};
  \node[font=\scriptsize\itshape, align=left, anchor=north west]
        at (-0.4,-1.1)
        {\;$\gamma \approx b_{\max}/d$ : the ratio gate is a parallax gate.\\
         \;Beyond $d_{\max}$ the linear system is ill-conditioned and\\
         \;the point is discarded before any $\chi^2$ test.};
\end{tikzpicture}
\caption[The baseline-ratio rejection]{Rejection is governed by parallax. With
a stereo baseline of $90$\,mm and $\beta=40$, every point beyond $3.6$\,m is
discarded while the vehicle is stationary.}
\label{fig:baselinegate}
\end{figure}

Evaluated for this platform, \eqref{eq:distceiling} is severe:

\begin{table}[htbp]
\centering
\begin{tabular}{@{}lrr@{}}
\toprule
condition & $b_{\max}$ (m) & $d_{\max}$ (m) \\
\midrule
at rest / pure pivot (stereo only)              & $0.090$ & $3.6$ \\
translating, $11$ clones at $0.2$\,m/s          & $0.240$ & $9.6$ \\
translating, $25$ clones at $0.2$\,m/s          & $0.420$ & $16.8$ \\
\bottomrule
\end{tabular}
\caption[Distance ceiling implied by the baseline ratio]{Indoor features
typically sit at $2$--$10$\,m. During a pure pivot --- the manoeuvre that
should be most informative about camera--IMU geometry --- almost the entire
scene is rejected.}
\label{tab:distceiling}
\end{table}

This mechanism is real and explains why enlarging the clone window helped
(peak features $2 \to 15$): it raises $b_{\max}$, hence $d_{\max}$. It is
\emph{not}, however, the binding constraint. Raising $\beta$ from $40$ to
$200$ made the measured outcome worse ($\bar{n}=0.75 \to 0.23$, peak
$25 \to 5$ over $777$ updates), and the sister estimator MINS already runs at
$\beta = 2000$ without benefit. Recorded here because the derivation is sound
and the ceiling is genuinely present --- it is simply not what closes the
channel.

% ─────────────────────────────────────────────────────────────────────────
\subsection{A structural blind spot: vertical unobservability in MINS}
\label{sec:minsz}

MINS survives where openVINS does not, because wheel odometry anchors it. The
wheel model is planar (\texttt{Wheel2DAng}); writing the measurement Jacobian
in the body frame,
\begin{equation}
\mathbf{h}_{\mathrm{wheel}}(\mathbf{x}) =
\begin{bmatrix} v_x \\ \omega_z \end{bmatrix},
\qquad
\frac{\partial \mathbf{h}}{\partial z} = \mathbf{0},
\label{eq:wheeljac}
\end{equation}
so the vertical coordinate receives \emph{no} information from the wheels.
Only the camera constrains $z$. When visual features are scarce --- the very
condition documented above --- the vertical channel is left to free inertial
integration and \eqref{eq:doubleintov} applies unchecked.

This was observed accidentally and then verified quantitatively. With the
robot stationary and the wheels reading zero, the horizontal estimate stayed
exact while $z$ diverged:
\begin{equation}
z(120\,\mathrm{s}) = -15.2\;\mathrm{m},
\qquad
x = 0.01\;\mathrm{m},
\qquad
y = -0.03\;\mathrm{m}.
\label{eq:minszdrift}
\end{equation}
Inverting \eqref{eq:doubleintov} recovers the driving residual,
\begin{equation}
\varepsilon_z = \frac{2\,|z|}{t^{2}}
              = \frac{2 \times 15.2}{120^{2}}
              = 2.1\times10^{-3}\;\mathrm{m/s^2},
\label{eq:epsz}
\end{equation}
against an independently measured specific-force residual of
$1.7\times10^{-3}$\,m/s$^2$ (Section~\ref{sec:accelexpire}). The two agree to
within $25\,\%$, which for a doubly-integrated quantity over two minutes is a
tight confirmation of the mechanism.

The practical consequence is a diagnostic asymmetry worth stating plainly:
\emph{a partially anchored estimator lies asymmetrically}. Inspecting $x$ and
$y$ --- the constrained axes --- gives full confidence while $z$ runs away
without limit and without any alarm. Every subsequent change to camera
extrinsics on this platform is now validated on the vertical axis as well,
through an explicit guard in the live supervisor.

% ─────────────────────────────────────────────────────────────────────────
\subsection{Seven hypotheses, seven measurements}
\label{sec:sevenhyp}

The residual failure --- $\bar{n}_{\mathrm{feat}} = 0.75$ against a capacity
of $40$ --- remains open. What is settled is what it is \emph{not}. Each row
below was tested on the robot and eliminated by measurement:

\begin{table}[htbp]
\centering
\small
\begin{tabular}{@{}p{47mm}p{52mm}r@{}}
\toprule
hypothesis & measured outcome & $\bar{n}$ \\
\midrule
scene lacks texture
  & MINS finds $50$ features in the same images
  & --- \\
global camera--IMU extrinsic
  & full Kalibr matrix performs \emph{worse} than the idealised one
  & $0.58$ \\
hybrid (Kalibr rotation, hand translation)
  & breaks the jointly-optimised pair; MINS falls $26$\,m in $z$
  & --- \\
$\chi^2$ gate too tight
  & $\times 5$ tolerance ($1 \to 2.2$\,px): $33\% \to 35\%$ non-empty
  & $0.75$ \\
translational parallax
  & $11 \to 25$ clones: peak $2 \to 15$, still starved
  & $0.75$ \\
stereo mis-declaration
  & \texttt{TrackKLT} derives epipolar geometry by RANSAC from data
  & --- \\
triangulation ratio $\beta$
  & $40 \to 200$: outcome degrades; MINS already at $2000$
  & $0.23$ \\
\bottomrule
\end{tabular}
\caption[Eliminated explanations for feature starvation]{Negative results,
recorded deliberately. Each cost an instrumented trial; repeating them would
cost the same again. The best configuration measured to date is
$25$ clones with a $\chi^2$ multiplier of $5$ on the idealised extrinsic:
$\bar{n} = 0.75$, peak $25$, $35\,\%$ non-empty updates.}
\label{tab:sevenhyp}
\end{table}

The instrumentation gap that remains is specific: nothing currently counts
stereo correspondences \emph{between} the matching stage and the triangulator.
Until that count exists, the two remaining candidate stages in
Figure~\ref{fig:featfunnel} cannot be separated, and any further parameter
change is tuning in the dark --- which this session has demonstrated does not
converge.
